\chapter{2009年上海卷（秋理）}

\section{填空题}

\begin{problem}
若复数 \(z\) 满足 \(z(1 + \i) = 1 - \i\)（\(\i\) 是虚数单位），则其共轭复数 \(\overline{z} =\)\fillinblank{}.
\end{problem}

\begin{answer}
\(\i\).
\end{answer}

\begin{solution}
由 \(z(1+\i)=1-\i\)，得
\[
z=\frac{1-\i}{1+\i}=\frac{(1-\i)^2}{(1+\i)(1-\i)}=\frac{-2\i}{2}=-\i.
\]
因此 \(\overline{z}=\i\).
\end{solution}

\begin{problem}
已知集合 \(A = \{x \mid x \leqslant 1\}\)，\(B = \{x \mid x \geqslant \alpha\}\)，且 \(A \cup B = \R\)，则实数 \(\alpha\) 的取值范围是\fillinblank{}.
\end{problem}

\begin{answer}
\(\alpha \leqslant 1\).
\end{answer}

\begin{solution}
由题意，\(A=(-\infty,1]\)，\(B=[\alpha,+\infty)\). 若 \(\alpha>1\)，则区间 \((1,\alpha)\) 中的实数既不属于 \(A\)，也不属于 \(B\)，从而 \(A\cup B\neq\R\). 反之，若 \(\alpha\leqslant1\)，则任意实数 \(x\leqslant1\) 时属于 \(A\)，任意实数 \(x>1\) 时有 \(x\geqslant\alpha\)，因而属于 \(B\)，所以 \(A\cup B=\R\). 故 \(\alpha\) 的取值范围为 \(\alpha\leqslant1\).
\end{solution}

\begin{problem}
若行列式 \(\begin{vmatrix} 4 & 5 & x \\ 1 & x & 3 \\ 7 & 8 & 9 \end{vmatrix}\) 中，元素 \(4\) 的代数余子式大于 \(0\)，则 \(x\) 满足的条件是\fillinblank{}.
\end{problem}

\begin{answer}
\(x > \frac{8}{3}\).
\end{answer}

\begin{solution}
元素 \(4\) 位于行列式的第 \(1\) 行第 \(1\) 列，因此它的代数余子式为
\[
C_{11}=(-1)^{1+1}\begin{vmatrix}x&3\\8&9\end{vmatrix}=9x-24.
\]
由题意 \(9x-24>0\)，所以 \(x>\dfrac{8}{3}\).
\end{solution}

\begin{problem}
某算法的程序框图如图所示，则输出量 \(y\) 与输入量 \(x\) 满足的关系式是\fillinblank{}.

\bitmapfigure[max width=0.36\linewidth]{img/2009/shanghai_science/q4_fig1.png}
\end{problem}

\begin{answer}
\[
y=\begin{cases}
2^x,&x\leqslant 1,\\
x-2,&x>1.
\end{cases}
\]
\end{answer}

\begin{solution}
由程序框图可知，当输入 \(x>1\) 时，沿“是”分支得到 \(y=x-2\)；当输入 \(x\leqslant 1\) 时，沿“否”分支得到 \(y=2^x\). 因此
\[
y=\begin{cases}
2^x,&x\leqslant 1,\\
x-2,&x>1.
\end{cases}
\]
\end{solution}

\begin{problem}
如图，若正四棱柱 \(ABCD - A_1B_1C_1D_1\) 的底面边长为 \(2\)，高为 \(4\)，则异面直线 \(BD_1\) 与 \(AD\) 所成角的大小是\fillinblank{}. （结果用反三角函数表示）

\bitmapfigure[max width=0.30\linewidth]{img/2009/shanghai_science/q5_fig1.png}
\end{problem}

\begin{answer}
\(\arctan\sqrt{5}\).
\end{answer}

\begin{solution}
建立直角坐标系，取 \(A=(0,0,0)\)，\(B=(2,0,0)\)，\(D=(0,2,0)\)，并令竖直方向为 \(z\) 轴，则 \(D_1=(0,2,4)\). 因而
\[
\overrightarrow{AD}=(0,2,0),\qquad \overrightarrow{BD_1}=(-2,2,4).
\]
设异面直线所成的锐角为 \(\theta\)，则
\[
\cos\theta=\frac{\left|\overrightarrow{AD}\cdot\overrightarrow{BD_1}\right|}{|\overrightarrow{AD}|\,|\overrightarrow{BD_1}|}=\frac{4}{2\sqrt{24}}=\frac{1}{\sqrt{6}}.
\]
于是 \(\tan\theta=\sqrt{\dfrac{1}{\cos^2\theta}-1}=\sqrt{5}\). 因为 \(0\leqslant\theta\leqslant\dfrac{\pi}{2}\)，所以 \(\theta=\arctan\sqrt{5}\).
\end{solution}

\begin{problem}
函数 \(y = 2\cos^2 x + \sin 2x\) 的最小值是\fillinblank{}.
\end{problem}

\begin{answer}
\(1 - \sqrt{2}\).
\end{answer}

\begin{solution}
利用 \(2\cos^2x=1+\cos2x\)，有
\[
y=1+\cos2x+\sin2x=1+\sqrt{2}\sin\left(2x+\frac{\pi}{4}\right).
\]
由于正弦函数的取值范围为 \([-1,1]\)，故
\[
y\geqslant1-\sqrt{2},
\]
且当 \(\sin\left(2x+\dfrac{\pi}{4}\right)=-1\) 时等号成立. 因此最小值为 \(1-\sqrt{2}\).
\end{solution}

\begin{problem}
某学校要从 \(5\) 名男生和 \(2\) 名女生中选出 \(2\) 人作为上海世博会志愿者，若用随机变量 \(\xi\) 表示选出的志愿者中女生的人数，则数学期望 \(E(\xi) =\)\fillinblank{}. （结果用最简分数表示）
\end{problem}

\begin{answer}
\(\frac{4}{7}\).
\end{answer}

\begin{solution}
从 \(5\) 名男生和 \(2\) 名女生中任取 \(2\) 人，共有 \(\mathrm C_7^2=21\) 种取法. 女生人数 \(\xi\) 的分布为
\[
P(\xi=0)=\frac{\mathrm C_5^2}{\mathrm C_7^2}=\frac{10}{21},\qquad P(\xi=1)=\frac{\mathrm C_5^1\mathrm C_2^1}{\mathrm C_7^2}=\frac{10}{21},
\]
\[
P(\xi=2)=\frac{\mathrm C_2^2}{\mathrm C_7^2}=\frac{1}{21}.
\]
因此
\[
E(\xi)=0\cdot\frac{10}{21}+1\cdot\frac{10}{21}+2\cdot\frac{1}{21}=\frac{4}{7}.
\]
\end{solution}

\begin{problem}
已知三个球的半径 \(R_1\)，\(R_2\)，\(R_3\) 满足 \(R_1 + 2R_2 = 3R_3\)，则它们的表面积 \(S_1\)，\(S_2\)，\(S_3\) 满足的等量关系是\fillinblank{}.
\end{problem}

\begin{answer}
\(\sqrt{S_1} + 2\sqrt{S_2} = 3\sqrt{S_3}\).
\end{answer}

\begin{solution}
球的表面积满足 \(S_i=4\pi R_i^2\). 由于半径为正数，故
\[
\sqrt{S_i}=2\sqrt{\pi}\,R_i\qquad(i=1,2,3).
\]
将题设关系 \(R_1+2R_2=3R_3\) 两边同乘 \(2\sqrt{\pi}\)，得到
\[
\sqrt{S_1}+2\sqrt{S_2}=3\sqrt{S_3}.
\]
\end{solution}

\begin{problem}
已知 \(F_1\)，\(F_2\) 是椭圆 \(C: \frac{x^2}{a^2} + \frac{y^2}{b^2} = 1\)（\(a > b > 0\)）的两个焦点，\(P\) 为椭圆 \(C\) 上一点，且 \(\overrightarrow{PF_1} \perp \overrightarrow{PF_2}\)，若 \(\triangle PF_1F_2\) 的面积为 \(9\)，则 \(b =\)\fillinblank{}.
\end{problem}

\begin{answer}
\(3\).
\end{answer}

\begin{solution}
设 \(PF_1=m\)，\(PF_2=n\)，焦半距为 \(c\). 由 \(\overrightarrow{PF_1}\perp\overrightarrow{PF_2}\)，三角形 \(PF_1F_2\) 在 \(P\) 处直角，故
\[
m^2+n^2=F_1F_2^2=(2c)^2.
\]
又由椭圆的定义，\(m+n=2a\). 两边平方并利用 \(a^2-b^2=c^2\)，得
\[
4a^2=m^2+n^2+2mn=4c^2+2mn,
\]
从而 \(mn=2(a^2-c^2)=2b^2\). 因此三角形面积为
\[
9=\frac12mn=b^2.
\]
由于 \(b>0\)，所以 \(b=3\).
\end{solution}

\begin{problem}
在极坐标系中，由三条直线 \(\theta = 0\)，\(\theta = \frac{\pi}{3}\)，\(\rho \cos \theta + \rho \sin \theta = 1\) 围成图形的面积是\fillinblank{}.
\end{problem}

\begin{answer}
\(\frac{3 - \sqrt{3}}{4}\).
\end{answer}

\begin{solution}
由直线 \(\rho\cos\theta+\rho\sin\theta=1\)，得
\[
\rho=\frac{1}{\cos\theta+\sin\theta}.
\]
所围图形对应于 \(0\leqslant\theta\leqslant\dfrac{\pi}{3}\)，故其面积为
\[
S=\frac12\int_0^{\pi/3}\frac{1}{(\cos\theta+\sin\theta)^2}\,\mathrm d\theta
=\frac14\int_0^{\pi/3}\sec^2\left(\theta-\frac{\pi}{4}\right)\,\mathrm d\theta
=\frac14\left[\tan\left(\theta-\frac{\pi}{4}\right)\right]_0^{\pi/3}.
\]
因为 \(\tan\dfrac{\pi}{12}=2-\sqrt3\)，所以
\[
S=\frac14\left((2-\sqrt3)-(-1)\right)=\frac{3-\sqrt3}{4}.
\]
\end{solution}

\begin{problem}
当 \(0 \leqslant x \leqslant 1\) 时，不等式 \(\sin \frac{\pi x}{2} \geqslant kx\) 成立. 则实数 \(k\) 的取值范围是\fillinblank{}.
\end{problem}

\begin{answer}
\(k \leqslant 1\).
\end{answer}

\begin{solution}
函数 \(g(x)=\sin\dfrac{\pi x}{2}\) 在 \([0,1]\) 上为凹函数，且 \(g(0)=0\)，\(g(1)=1\). 由凹函数图象位于端点连线的上方，得
\[
\sin\frac{\pi x}{2}\geqslant x\qquad(0\leqslant x\leqslant1).
\]
若题设不等式对所有 \(x\in[0,1]\) 成立，取 \(x=1\) 得 \(1\geqslant k\)，即 \(k\leqslant1\). 反之，当 \(k\leqslant1\) 时，由 \(x\geqslant0\) 有 \(kx\leqslant x\leqslant\sin\dfrac{\pi x}{2}\)，题设不等式对整个区间成立. 因此 \(k\leqslant1\).
\end{solution}

\begin{problem}
已知函数 \(f(x) = \sin x + \tan x\)，项数为 \(27\) 的等差数列 \(\{a_n\}\) 满足 \(a_n \in \left( -\frac{\pi}{2}, \frac{\pi}{2} \right)\)，且公差 \(d \neq 0\). 若 \(f(a_1) + f(a_2) + \cdots + f(a_{27}) = 0\)，则当 \(k =\)\fillinblank{} 时，\(f(a_k) = 0\).
\end{problem}

\begin{answer}
\(14\).
\end{answer}

\begin{solution}
函数 \(f(x)=\sin x+\tan x\) 在 \(\left(-\dfrac{\pi}{2},\dfrac{\pi}{2}\right)\) 上为奇函数，且
\[
f'(x)=\cos x+\sec^2x>0,
\]
所以 \(f\) 在该区间上严格递增. 等差数列的中项为 \(a_{14}\)，并且第 \(14-j\) 项和第 \(14+j\) 项可写成 \(a_{14}-u\) 与 \(a_{14}+u\)，其中 \(u=|jd|>0\).

若 \(a_{14}>0\)，当 \(a_{14}-u\geqslant0\) 时，两个函数值之和为正；当 \(a_{14}-u<0\) 时，由严格递增性和奇函数性质，
\[
a_{14}+u>u-a_{14}>0\quad\Longrightarrow\quad f(a_{14}+u)>f(u-a_{14})=-f(a_{14}-u),
\]
故仍有 \(f(a_{14}+u)+f(a_{14}-u)>0\). 所有成对项的和以及中项 \(f(a_{14})\) 都为正，于是全体和不可能为零. 同理，若 \(a_{14}<0\)，全体和为负.

因此题设等式只能在 \(a_{14}=0\) 时成立. 此时 \(f(a_{14})=f(0)=0\)，所以 \(k=14\).
\end{solution}

\begin{problem}
某地街道呈现东—西，南—北向的网格状，相邻街距都为 \(1\). 两街道相交的点称为格点. 若以互相垂直的两条街道为轴建立直角坐标系，现有下述格点 \((-2, 2)\)，\((3, 1)\)，\((3, 4)\)，\((-2, 3)\)，\((4, 5)\)，\((6, 6)\) 为报刊零售点. 请确定一个格点（除零售点外）为发行站，使 \(6\) 个零售点沿街道到发行站之间路程的和最短\fillinblank{}.
\end{problem}

\begin{answer}
\((3, 3)\).
\end{answer}

\begin{solution}
设发行站为格点 \((u,v)\). 六个零售点到它的沿街路程之和为
\[
\begin{aligned}
T(u,v)&=|u+2|+|u-3|+|u-3|+|u+2|+|u-4|+|u-6|\\
&\quad+|v-2|+|v-1|+|v-4|+|v-3|+|v-5|+|v-6|.
\end{aligned}
\]
横、纵坐标部分彼此独立. 横坐标数据按从小到大为 \(-2,-2,3,3,4,6\)，其和绝对值之和在中位数区间 \([3,3]\) 处取得最小值，所以 \(u=3\). 纵坐标数据按从小到大为 \(1,2,3,4,5,6\)，其和绝对值之和在区间 \([3,4]\) 内取得最小值，所以 \(v\) 可取 \(3\) 或 \(4\). 点 \((3,4)\) 是零售点，不能作为发行站，故只能取 \((3,3)\).
\end{solution}

\begin{problem}
将函数 \(y = \sqrt{4 + 6x - x^2} - 2\)（\(x \in [0, 6]\)）的图象绕坐标原点逆时针方向旋转角 \(\theta\)（\(0 \leqslant \theta \leqslant \alpha\)），得到曲线 \(C\). 若对于每一个旋转角 \(\theta\)，曲线 \(C\) 都是一个函数的图象，则 \(\alpha\) 的最大值为\fillinblank{}.
\end{problem}

\begin{answer}
\(\arctan\frac{2}{3}\).
\end{answer}

\begin{solution}
原曲线可参数化为 \((x,y(x))\)，其中
\[
y(x)=\sqrt{13-(x-3)^2}-2,\qquad 0\leqslant x\leqslant6.
\]
旋转 \(\theta\) 后点的横坐标为
\[
X_\theta(x)=x\cos\theta-y(x)\sin\theta.
\]
若旋转后的曲线是函数图象，横坐标必须沿曲线单调，从而需考察
\[
X_\theta'(x)=\cos\theta-y'(x)\sin\theta,\qquad y'(x)=\frac{3-x}{\sqrt{13-(x-3)^2}}.
\]
在 \([0,6]\) 上，\(y'(x)\leqslant y'(0)=\dfrac32\)，因此当 \(0\leqslant\theta\leqslant\arctan\dfrac23\) 时，
\[
X_\theta'(x)\geqslant\cos\theta-\frac32\sin\theta\geqslant0.
\]
除端点情形外导数为正，故 \(X_\theta(x)\) 严格递增，旋转后的曲线是函数图象.

若 \(\theta>\arctan\dfrac23\) 且 \(\theta<\dfrac\pi2\)，则
\[
X_\theta'(0)=\cos\theta-\frac32\sin\theta<0,\qquad X_\theta'(6)=\cos\theta+\frac32\sin\theta>0.
\]
又因
\[
y''(x)=-\frac{13}{\left(13-(x-3)^2\right)^{3/2}}<0,
\]
故 \(X_\theta'(x)\) 在 \([0,6]\) 上严格递增，横坐标函数先减后增，在某个内点取得唯一最小值. 取最小值附近同一水平线上的两个不同点，旋转后它们的横坐标相同而点本身不同，因而纵坐标不同，旋转后的曲线不是函数图象. 因而允许的最大旋转角为 \(\arctan\dfrac23\).
\end{solution}

\section{单选题}

\begin{problem}
“\(-2 \leqslant a \leqslant 2\)”是“实系数一元二次方程 \(x^2 + ax + 1 = 0\) 有虚根”的（\quad）

\choices
    {必要不充分条件}
    {充分不必要条件}
    {充要条件}
    {既不充分也不必要条件}
\end{problem}

\begin{answer}
A.
\end{answer}

\begin{solution}
方程 \(x^2+ax+1=0\) 有虚根的充要条件是判别式小于零，即
\[
\Delta=a^2-4<0\quad\Longleftrightarrow\quad -2<a<2.
\]
因此 \(-2\leqslant a\leqslant2\) 是该条件的必要条件，但在 \(a=2\) 或 \(a=-2\) 时方程有重根而不是虚根，故不是充分条件. 所以应选 A.
\end{solution}

\begin{problem}
若事件 \(E\) 与 \(F\) 相互独立，且 \(P(E) = P(F) = \frac{1}{4}\)，则 \(P(E \cap F)\) 的值等于（\quad）

\choices
    {\(0\)}
    {\(\frac{1}{16}\)}
    {\(\frac{1}{4}\)}
    {\(\frac{1}{2}\)}
\end{problem}

\begin{answer}
B.
\end{answer}

\begin{solution}
由事件 \(E\) 与 \(F\) 相互独立，
\[
P(E\cap F)=P(E)P(F)=\frac14\cdot\frac14=\frac1{16}.
\]
所以应选 B.
\end{solution}

\begin{problem}
在发生某公共卫生事件期间，有专业机构认为该事件在一段时间没有发生规模群体感染的标志为“连续 \(10\) 天，每天新增疑似病例不超过 \(7\) 人”. 根据过去 \(10\) 天甲，乙，丙，丁四地新增疑似病例数据，一定符合该标志的是（\quad）

\choices
    {甲地：总体均值为 \(3\)，中位数为 \(4\)}
    {乙地：总体均值为 \(1\)，总体方差大于 \(0\)}
    {丙地：中位数为 \(2\)，众数为 \(3\)}
    {丁地：总体均值为 \(2\)，总体方差为 \(3\)}
\end{problem}

\begin{answer}
D.
\end{answer}

\begin{solution}
设过去十天的新增疑似病例数为非负整数 \(x_1,\ldots,x_{10}\). 对丁地，均值为 \(2\)，方差为 \(3\)，所以
\[
\sum_{i=1}^{10}x_i=20,\qquad \sum_{i=1}^{10}x_i^2=10(3+2^2)=70.
\]
若某一天 \(x_i=t\geqslant8\)，则其余九天的和为 \(20-t\). 由均方不等式，
\[
\sum_{j\ne i}x_j^2\geqslant\frac{(20-t)^2}{9},
\]
从而
\[
\sum_{j=1}^{10}x_j^2\geqslant t^2+\frac{(20-t)^2}{9}\geqslant 8^2+\frac{12^2}{9}=80>70,
\]
矛盾. 因此每天新增病例数都不超过 \(7\)，丁地一定符合标志.

其余选项均不能保证最大值不超过 \(7\). 例如甲地可取数据 \(0,0,0,0,4,4,4,4,4,10\)，乙地可取 \(0,0,0,0,0,0,0,0,0,10\)，丙地可取 \(0,0,0,1,1,3,3,3,3,10\)，它们分别满足相应条件但都出现了超过 \(7\) 的数据. 所以应选 D.
\end{solution}

\begin{problem}
过圆 \(C: (x - 1)^2 + (y - 1)^2 = 1\) 的圆心作直线分别交 \(x\)，\(y\) 正半轴于点 \(A\)，\(B\)，\(\triangle AOB\) 被圆分成四部分（如图），若这四部分图形面积满足 \(S_{\mathrm I} + S_{\mathrm{IV}} = S_{\mathrm{II}} + S_{\mathrm{III}}\)，则直线 \(AB\) 有（\quad）

\bitmapfigure[max width=0.42\linewidth]{img/2009/shanghai_science/q18_fig1.png}

\choices
    {\(0\) 条}
    {\(1\) 条}
    {\(2\) 条}
    {\(3\) 条}
\end{problem}

\begin{answer}
B.
\end{answer}

\begin{solution}
设直线与两坐标轴的交点为 \(A=(a,0)\)，\(B=(0,b)\). 因为直线过圆心 \(C=(1,1)\)，所以
\[
\frac1a+\frac1b=1.
\]
令 \(t=a-1>0\)，则 \(b-1=\dfrac1t\)，且每个正数 \(t\) 唯一确定一条这样的直线.

区域 \(S_{\mathrm{II}}\) 由两坐标轴和圆在第一象限的弧围成，其面积是单位正方形面积减去四分之一单位圆面积，即
\[
S_{\mathrm{II}}=1-\frac\pi4.
\]
直线 \(AB\) 过圆心，且三角形 \(AOB\) 位于直线 \(AB\) 的下侧，所以圆在三角形 \(AOB\) 内的部分 \(S_{\mathrm{IV}}\) 是一个半圆，从而
\[
S_{\mathrm{IV}}=\frac\pi2.
\]

令 \(T_A=(1,0)\)，并令 \(\phi=\arctan t\). 设直线与圆在靠近 \(A\) 的交点为 \(P\)，则由方向向量 \((t,-1)\) 得
\[
P=C+\frac{(t,-1)}{\sqrt{1+t^2}}.
\]
三角形 \(AT_AP\) 的面积为 \(\dfrac12t\left(1-\dfrac1{\sqrt{1+t^2}}\right)\)，弦 \(T_AP\) 与弧 \(T_AP\) 所夹的圆弓形面积为 \(\dfrac12\left(\phi-\dfrac{t}{\sqrt{1+t^2}}\right)\)，因此
\[
S_{\mathrm I}=\frac{t-\phi}{2}=\frac{t-\arctan t}{2}.
\]
交换 \(x\)，\(y\) 后，区域 \(S_{\mathrm{III}}\) 的对应参数为 \(1/t\)，同样计算得
\[
S_{\mathrm{III}}=\frac{t^{-1}-\arctan(t^{-1})}{2}.
\]
于是原条件等价于
\[
S_{\mathrm I}-S_{\mathrm{III}}=S_{\mathrm{II}}-S_{\mathrm{IV}}=1-\frac{3\pi}{4}.
\]
令
\[
G(t)=S_{\mathrm I}-S_{\mathrm{III}}=\frac12\left(t-\frac1t-\arctan t+\arctan\frac1t\right).
\]
则
\[
G'(t)=\frac12\left(1+\frac1{t^2}-\frac2{1+t^2}\right)=\frac{t^4+1}{2t^2(1+t^2)}>0\qquad(t>0).
\]
并且 \(G(t)\to-\infty\)（\(t\to0^+\)），\(G(t)\to+\infty\)（\(t\to+\infty\)），所以方程 \(G(t)=1-\dfrac{3\pi}{4}\) 恰有一个正数解，即恰有一条直线满足条件. 所以应选 B.
\end{solution}

\section{解答题}

\begin{problem}
如图，在直三棱柱 \(ABC - A_1B_1C_1\) 中，\(AA_1 = BC = AB = 2\)，\(AB \perp BC\)，求二面角 \(B_1 - A_1C - C_1\) 的大小.

\bitmapfigure[max width=0.45\linewidth]{img/2009/shanghai_science/q19_fig1.png}
\end{problem}

\begin{solution}
建立空间直角坐标系，取 \(B\) 为原点，\(BA\)，\(BC\)，\(BB_1\) 所在直线分别为 \(x\) 轴、\(y\) 轴、\(z\) 轴，则
\[
B=(0,0,0),\quad A=(2,0,0),\quad C=(0,2,0),
\]
\[
A_1=(2,0,2),\quad B_1=(0,0,2),\quad C_1=(0,2,2).
\]
设 \(\bs{u}=\overrightarrow{A_1C}=(-2,2,-2)\). 在平面 \(B_1A_1C\) 内，取从棱 \(A_1C\) 指向 \(B_1\) 且垂直于棱 \(A_1C\) 的向量
\[
\bs{v}_1=\overrightarrow{A_1B_1}-\frac{\overrightarrow{A_1B_1}\cdot\bs{u}}{\bs{u}\cdot\bs{u}}\bs{u}=(-2,0,0)-\frac{4}{12}(-2,2,-2)=\frac{2}{3}(-2,-1,1).
\]
在平面 \(A_1CC_1\) 内，取从棱 \(A_1C\) 指向 \(C_1\) 且垂直于棱 \(A_1C\) 的向量
\[
\bs{v}_2=\overrightarrow{A_1C_1}-\frac{\overrightarrow{A_1C_1}\cdot\bs{u}}{\bs{u}\cdot\bs{u}}\bs{u}=(-2,2,0)-\frac{8}{12}(-2,2,-2)=\frac{2}{3}(-1,1,2).
\]
所以所求二面角等于 \(\bs{v}_1\) 与 \(\bs{v}_2\) 的夹角 \(\theta\)，从而
\[
\cos\theta=\frac{(-2,-1,1)\cdot(-1,1,2)}{\sqrt{(-2)^2+(-1)^2+1^2}\sqrt{(-1)^2+1^2+2^2}}=\frac{1}{2}.
\]
由于二面角的范围是 \([0,\pi]\)，故 \(\theta=\frac{\pi}{3}\).
\end{solution}

\begin{problem}
有时可用函数 \(f(x) = \begin{cases}
0.1 + 15 \ln \frac{a}{a - x}, & x \leqslant 6,\\
\frac{x - 4.4}{x - 4}, & x > 6,
\end{cases}\) 描述学习某学科知识的掌握程度，其中 \(x\) 表示某学科知识的学习次数（\(x \in \mathbf{N}^*\)），\(f(x)\) 表示对该学科知识的掌握程度，正实数 \(a\) 与学科知识有关.

\begin{enumerate}
\item 证明：当 \(x \geqslant 7\) 时，掌握程度的增加量 \(f(x + 1) - f(x)\) 总是下降；

\item 根据经验，学科甲，乙，丙对应的 \(a\) 的取值范围分别为 \((115, 121]\)，\((121, 127]\)，\((121, 133]\). 当学习某学科知识 \(6\) 次时，掌握程度是 \(85\%\)，请确定相应的学科.
\end{enumerate}
\end{problem}

\begin{solution}
\begin{enumerate}
\item 当 \(x\geqslant 7\) 时，\(x\) 和 \(x+1\) 都大于 \(6\)，因此
\[
f(x)=\frac{x-4.4}{x-4}=1-\frac{0.4}{x-4}.
\]
令 \(\Delta_x=f(x+1)-f(x)\)，则
\[
\Delta_x=\frac{0.4}{(x-4)(x-3)}.
\]
于是
\[
\Delta_{x+1}-\Delta_x=\frac{0.4}{(x-3)(x-2)}-\frac{0.4}{(x-4)(x-3)}=-\frac{0.8}{(x-4)(x-3)(x-2)}<0.
\]
所以当 \(x\geqslant 7\) 时，掌握程度的增加量总是下降.

\item 由题意，\(f(6)=0.85\)，且 \(6\leqslant 6\)，故
\[
0.1+15\ln\frac{a}{a-6}=0.85,
\]
即
\[
\ln\frac{a}{a-6}=\frac{1}{20},\qquad \frac{a}{a-6}=\e^{1/20}.
\]
解得
\[
a=\frac{6\e^{1/20}}{\e^{1/20}-1}=\frac{6}{1-\e^{-1/20}}\approx123.025.
\]
因此 \(a\in(121,127]\cap(121,133]\)，相应的学科为乙、丙.
\end{enumerate}
\end{solution}

\begin{problem}
已知双曲线 \(C: \frac{x^2}{2} - y^2 = 1\)，设过点 \(A(-3\sqrt{2}, 0)\) 的直线 \(l\) 的方向向量 \(\bs{e} = (1, k)\).

\begin{enumerate}
\item 当直线 \(l\) 与双曲线 \(C\) 的一条渐近线 \(m\) 平行时，求直线 \(l\) 的方程及 \(l\) 与 \(m\) 的距离；

\item 证明：当 \(k > \frac{\sqrt{2}}{2}\) 时，在双曲线 \(C\) 的右支上不存在点 \(Q\)，使之到直线 \(l\) 的距离为 \(\sqrt{6}\).
\end{enumerate}
\end{problem}

\begin{solution}
\begin{enumerate}
\item 直线 \(l\) 过点 \(A(-3\sqrt{2},0)\)，方向向量为 \(\bs{e}=(1,k)\)，故
\[
l:y=k(x+3\sqrt{2}),\qquad kx-y+3\sqrt{2}k=0.
\]
双曲线 \(C\) 的渐近线为
\[
m_1:y=\frac{x}{\sqrt{2}},\qquad m_2:y=-\frac{x}{\sqrt{2}}.
\]
当 \(k=\frac{\sqrt{2}}{2}\) 时，\(l\parallel m_1\)，且
\[
l:y=\frac{x}{\sqrt{2}}+3,
\]
所以
\[
d(l,m_1)=\frac{3\sqrt{2}}{\sqrt{1+2}}=\sqrt{6}.
\]
当 \(k=-\frac{\sqrt{2}}{2}\) 时，\(l\parallel m_2\)，且
\[
l:y=-\frac{x}{\sqrt{2}}-3,
\]
同理有 \(d(l,m_2)=\sqrt{6}\). 因此，直线 \(l\) 的方程及其与相应渐近线的距离分别为
\[
k=\frac{\sqrt{2}}{2}:\quad l:y=\frac{x}{\sqrt{2}}+3,\quad d=\sqrt{6};
\]
\[
k=-\frac{\sqrt{2}}{2}:\quad l:y=-\frac{x}{\sqrt{2}}-3,\quad d=\sqrt{6}.
\]

\item 设 \(Q(x,y)\) 在双曲线 \(C\) 的右支上，则 \(x\geqslant\sqrt{2}\)，且
\[
y=\pm\sqrt{\frac{x^2}{2}-1}.
\]
点 \(Q\) 到直线 \(l\) 的距离为
\[
d(Q,l)=\frac{\left|kx-y+3\sqrt{2}k\right|}{\sqrt{k^2+1}}.
\]
先考虑双曲线的下支. 此时
\[
kx-y+3\sqrt{2}k=kx+\sqrt{\frac{x^2}{2}-1}+3\sqrt{2}k\geqslant4\sqrt{2}k.
\]
由 \(k>\frac{\sqrt{2}}{2}\) 得
\[
(4\sqrt{2}k)^2-6(k^2+1)=26k^2-6>0,
\]
所以 \(kx-y+3\sqrt{2}k>\sqrt{6(k^2+1)}\)，从而 \(d(Q,l)>\sqrt{6}\).

再考虑双曲线的上支. 令
\[
g(x)=kx-\sqrt{\frac{x^2}{2}-1},\qquad x\geqslant\sqrt{2}.
\]
因为 \(k>\frac{\sqrt{2}}{2}\)，且
\[
g'(x)=k-\frac{x}{2\sqrt{x^2/2-1}}.
\]
令 \(g'(x)=0\)，得
\[
x_0=\frac{2k}{\sqrt{2k^2-1}}
\]
；由 \(g'(x)\) 在 \(x_0\) 左侧为负、右侧为正，知 \(g(x)\) 在 \(x_0\) 处取得最小值，且
\[
g(x_0)=\sqrt{2k^2-1}.
\]
于是
\[
kx-y+3\sqrt{2}k=g(x)+3\sqrt{2}k\geqslant3\sqrt{2}k+\sqrt{2k^2-1}>0.
\]
并且
\[
\left(3\sqrt{2}k+\sqrt{2k^2-1}\right)^2-6(k^2+1)
=7(2k^2-1)+6\sqrt{2}k\sqrt{2k^2-1}>0.
\]
所以仍有 \(d(Q,l)>\sqrt{6}\). 上、下两支均不可能出现距离等于 \(\sqrt{6}\) 的点，命题得证.
\end{enumerate}
\end{solution}

\begin{problem}
已知函数 \(y = f^{-1}(x)\) 是 \(y = f(x)\) 的反函数. 定义：若对给定的实数 \(a\)（\(a \neq 0\)），函数 \(y = f(x + a)\) 与 \(y = f^{-1}(x + a)\) 互为反函数，则称 \(y = f(x)\) 满足“\(a\) 和性质”；若函数 \(y = f(ax)\) 与 \(y = f^{-1}(ax)\) 互为反函数，则称 \(y = f(x)\) 满足“\(a\) 积性质”.

\begin{enumerate}
\item 判断函数 \(g(x) = x^2\)（\(x > 0\)）是否满足“\(1\) 和性质”，并说明理由；

\item 求所有满足“\(2\) 和性质”的一次函数；

\item 设函数 \(y = f(x)\)（\(x > 0\)）对任何 \(a > 0\) 都满足“\(a\) 积性质”，求 \(y = f(x)\) 的表达式.
\end{enumerate}
\end{problem}

\begin{solution}
\begin{enumerate}
\item 当 \(x>0\) 时，\(g^{-1}(x)=\sqrt{x}\). 若记
\[
u(x)=g(x+1)=(x+1)^2,\qquad v(x)=g^{-1}(x+1)=\sqrt{x+1},
\]
则
\[
u\bigl(v(x)\bigr)=\left(\sqrt{x+1}+1\right)^2\neq x.
\]
因此 \(u(x)\) 与 \(v(x)\) 不是互为反函数，函数 \(g(x)=x^2\) 不满足“\(1\) 和性质”.

\item 设一次函数 \(f(x)=rx+s\)，其中 \(r\neq0\)，则
\[
f^{-1}(x)=\frac{x-s}{r}.
\]
令
\[
u(x)=f(x+2)=rx+2r+s,\qquad v(x)=f^{-1}(x+2)=\frac{x+2-s}{r}.
\]
若 \(u\) 与 \(v\) 互为反函数，则 \(u(v(x))=x\)，而
\[
u\bigl(v(x)\bigr)=x+2+2r.
\]
故 \(r=-1\). 反之，当 \(r=-1\) 时，\(u(x)=v(x)=-x-2+s\)，且
\[
u\bigl(u(x)\bigr)=x,
\]
所以所有满足“\(2\) 和性质”的一次函数为
\[
y=-x+s\qquad(s\in\R).
\]

\item 记 \(h=f^{-1}\). 对任意 \(a>0\)，函数 \(x\mapsto f(ax)\) 与 \(x\mapsto h(ax)\) 互为反函数，因此
\[
f\bigl(a h(ax)\bigr)=x.
\]
对等式两边施加 \(h\)，得到
\[
a h(ax)=h(x),\qquad h(ax)=\frac{h(x)}{a}.
\]
任取正数 \(t,x\)，令 \(a=t/x\)，则
\[
h(t)=\frac{x}{t}h(x),\qquad th(t)=xh(x).
\]
所以存在常数 \(c\)，使得
\[
h(t)=\frac{c}{t}.
\]
由于 \(h(t)\) 的值属于函数 \(f\) 的定义域 \((0,+\infty)\)，故 \(c>0\). 于是
\[
f(x)=\frac{c}{x}\qquad(x>0).
\]
反之，若 \(f(x)=c/x\)，其中 \(c>0\)，则 \(f^{-1}(x)=c/x\)，且对任意 \(a>0\)，
\[
f(ax)=f^{-1}(ax)=\frac{c}{ax},
\]
并且
\[
f\left(a\cdot\frac{c}{ax}\right)=x.
\]
所以该函数确实对任意正数 \(a\) 满足“\(a\) 积性质”. 因此
\[
y=\frac{c}{x}\qquad(c>0).
\]
\end{enumerate}
\end{solution}

\begin{problem}
已知 \(\{a_n\}\) 是公差为 \(d\) 的等差数列，\(\{b_n\}\) 是公比为 \(q\) 的等比数列.

\begin{enumerate}
\item 若 \(a_n = 3n + 1\)，是否存在 \(m\)，\(k \in \mathbf{N}^*\)，有 \(a_m + a_{m+1} = a_k\)？请说明理由；

\item 找出所有数列 \(\{a_n\}\) 和 \(\{b_n\}\)，使对一切 \(n \in \mathbf{N}^*\)，\(\frac{a_{n+1}}{a_n} = b_n\)，并说明理由；

\item 若 \(a_1 = 5\)，\(d = 4\)，\(b_1 = q = 3\)，试确定所有的 \(p\)，使数列 \(\{a_n\}\) 中存在某个连续 \(p\) 项的和是数列 \(\{b_n\}\) 中的一项，请证明.
\end{enumerate}
\end{problem}

\begin{solution}
\begin{enumerate}
\item 因为 \(a_n=3n+1\)，所以
\[
a_m+a_{m+1}=(3m+1)+(3m+4)=6m+5.
\]
若存在正整数 \(m,k\) 使 \(a_m+a_{m+1}=a_k\)，则
\[
6m+5=3k+1,
\]
即 \(3(k-2m)=4\)，这不可能. 因此不存在这样的 \(m,k\).

\item 由于题设中的比值对一切正整数 \(n\) 有意义，所以 \(a_n\neq0\)，且 \(b_n\neq0\). 又因为 \(b_{n+1}=qb_n\)，故
\[
q=\frac{b_{n+1}}{b_n}=\frac{a_{n+2}/a_{n+1}}{a_{n+1}/a_n}=\frac{a_na_{n+2}}{a_{n+1}^2}.
\]
等差数列满足 \(a_n=a_{n+1}-d\)，\(a_{n+2}=a_{n+1}+d\)，所以
\[
q=\frac{(a_{n+1}-d)(a_{n+1}+d)}{a_{n+1}^2}=1-\frac{d^2}{a_{n+1}^2}.
\]
若 \(d\neq0\)，则 \(a_{n+1}^2\) 对所有正整数 \(n\) 都相同. 于是相邻三项满足
\[
a_{n+1}^2=a_{n+2}^2=a_{n+3}^2.
\]
由 \(a_{n+2}-a_{n+1}=a_{n+3}-a_{n+2}=d\neq0\)，前两个等式分别迫使
\[
a_{n+2}=-a_{n+1},\qquad a_{n+3}=-a_{n+2}=a_{n+1},
\]
从而 \(a_{n+3}-a_{n+1}=2d=0\)，矛盾. 所以 \(d=0\).

此时 \(a_n=a_1\neq0\)，由题设 \(b_n=a_{n+1}/a_n=1\)，故 \(b_1=1\)，\(q=1\). 反之，任取非零常数 \(c\)，令 \(a_n=c\)，\(b_n=1\)，即可满足题设条件. 因此所有数列为
\[
a_n=c\quad(c\neq0),\qquad b_n=1.
\]

\item 由 \(a_1=5\)，\(d=4\)，得
\[
a_n=5+4(n-1)=4n+1.
\]
又因 \(b_1=q=3\)，所以 \(b_n=3^n\). 从第 \(m\) 项起的连续 \(p\) 项和为
\[
\sum_{j=0}^{p-1}a_{m+j}=\sum_{j=0}^{p-1}(4m+4j+1)=p(4m+1)+4\cdot\frac{p(p-1)}{2}=p(4m+2p-1).
\]
若该和等于某一项 \(b_n=3^n\)，则 \(p\) 是 \(3^n\) 的因数，故必有 \(p=3^r\)，其中 \(r=0,1,2,\ldots\). 设
\[
4m+2\cdot3^r-1=3^s.
\]
由于 \(m\geqslant1\)，左端大于 \(1\)，且
\[
3^s\equiv2\cdot3^r-1\equiv1\pmod 4,
\]
所以 \(s\) 为偶数. 这说明上述必要条件没有进一步限制 \(r\).

反之，任取 \(r=0,1,2,\ldots\)，取 \(s=2r+2\)，并令
\[
m=\frac{3^{2r+2}-2\cdot3^r+1}{4}.
\]
由于 \(3^{2r+2}\equiv1\pmod4\)，且 \(2\cdot3^r-1\equiv1\pmod4\)，故 \(m\) 是正整数，并且
\[
p(4m+2p-1)=3^r\cdot3^{2r+2}=3^{3r+2}=b_{3r+2}
\]
当 \(p=3^r\) 时成立. 因此所有可能的 \(p\) 为
\[
p=3^r\qquad(r=0,1,2,\ldots).
\]
\end{enumerate}
\end{solution}
